% Part: first-order-logic
% Chapter: semantics
% Section: intro-semantics

\documentclass[../../../include/open-logic-section]{subfiles}

\begin{document}

\olfileid{fol}{syn}{its}

\olsection{Pambuka}

Menehi teges marang ungkapan iku wewengkone semantik. Gagasan pokok
ing semantik yaiku nyukupi ing !!a{structure}. !!^a{structure} menehi
teges marang unsur-unsur dhasar basa: !!a{domain} iku himpunan obyek
kang ora kosong. Kuantor ditafsirake kanthi variabele njupuk nilai saka domain iki,
!!{constant}s diwenehi elemen ing domain, !!{function}s diwenehi
fungsi kang njupuk tupel unsur saka !!{domain} miturut aritase lan ngasilake unsur ing domain iku, lan !!{predicate}s diwenehi
relasi ing !!{domain}. !!{domain} bebarengan karo pangaji marang
kosakata dhasar mbentuk !!a{structure}. !!^{variable}s bisa muncul
ing !!{formula}s, lan supaya bisa menehi semantik, kita uga kudu
menehi !!{element}s saka !!{domain} marang variabel kasebut---iki
diarani pangaji variabel. Pungkasan, relasi nyukupi nyawijikake
kabeh iki. !!^a{formula} bisa dicukupi ing !!a{structure}~$\Struct{M}$
relatif marang pangaji !!a{variable}~$s$, kang ditulis
$\Sat{M}{!A}[s]$. Relasi iki uga ditetepake kanthi induksi marang
struktur~$!A$, nggunakake tabel bebener kanggo konektif logis
supaya bisa nemtokake, upamane, kapan $(!A \land !B)$ dicukupi
adhedhasar dicukupi utawa orane $!A$ lan ~$!B$. Banjur kabukten
yen pangaji !!{variable} ora nduweni pangaribawa nalika
!!{formula}~$!A$ iku !!a{sentence}, yaiku ora nduweni variabel bebas,
mula kita bisa kandha yen !!{sentence}s mung dicukupi (utawa ora)
ing !!{structure}s. Cathetan koreksi sumber OLPL-748: fungsi n-panggonan njupuk tupel n unsur domain lan ngasilake siji unsur domain, kaya definisi formal ing bagean sabanjure; frasa sumber ``from the domain to itself'' mung trep nalika n=1.

Adhedhasar relasi nyukupi $\Sat{M}{!A}$ kanggo !!{sentence}s, kita
banjur bisa nemtokake gagasan semantis dhasar: validitas, akibat
semantis, lan bisa dicukupi. !!^a{sentence} iku valid, $\Entails !A$,
yen saben !!{structure} nyukupi ukara kasebut. Ukara iku dadi akibat
saka sawijining himpunan !!{sentence}s, $\Gamma \Entails !A$,
yen saben !!{structure} kang nyukupi kabeh !!{sentence}s ing~$\Gamma$
uga nyukupi~$!A$. Lan sawijining himpunan !!{sentence}s bisa
dicukupi yen ana sawenehing !!{structure} kang nyukupi kabeh
!!{sentence}s ing himpunan iku bebarengan. Amarga !!{formula}s
ditetepake kanthi induktif, lan relasi nyukupi sabanjure uga
ditetepake kanthi induksi marang struktur !!{formula}s, kita
bisa nggunakake induksi kanggo mbuktekake sipat-sipat semantik
kita lan ngubungake gagasan-gagasan semantis kang wis ditetepake.

\end{document}
